
\documentclass[12pt,a4paper]{article}

\usepackage[brazil]{babel}
\usepackage{fontspec}
\usepackage{xeCJK}
\usepackage[table,xcdraw]{xcolor}
\usepackage{ruby}
\usepackage{amsmath}
\usepackage{float}

\setmainfont{Latin Modern Roman}
\setCJKmainfont{Noto Serif CJK JP}

\begin{document}

\section*{Aulas de Língua Japonesa Seishun}

\subsection*{Informações Gerais}

\begin{table}[H]
\begin{tabular}{|l|l|lll}
\cline{1-2}
Período de Realização    & De 3 de setembro a 26 de novembro de 2026               &  &  &  \\ \cline{1-2}
Carga Horária Total      & 26 horas (13 encontros de 2 horas)                      &  &  &  \\ \cline{1-2}
Horário                  & Quintas-feiras, das 19h às 21h                          &  &  &  \\ \cline{1-2}
Local                    & UTFPR Sala (definir)                                          &  &  &  \\ \cline{1-2}
Quantidade de Vagas      & 40                                                      &  &  &  \\ \cline{1-2}
Custo                    & R\$ 0,00                                                &  &  &  \\ \cline{1-2}
Público-alvo             & Pessoas com idade entre 14 e 50 anos                    &  &  &  \\ \cline{1-2}
Perfil dos Participantes & Pessoas interessadas em aprender a língua japonesa      &  &  &  \\ \cline{1-2}
Pré-requisitos           & Nenhum                                                  &  &  &  \\ \cline{1-2}
Divulgação               & Comunidade interna e externa, por meio das redes sociais &  &  &  \\ \cline{1-2}
\end{tabular}
\end{table}

\subsection*{Objetivo Geral}

Realizar aulas de japonês básico, apresentando aos participantes aspectos do idioma e da cultura japonesa. As aulas terão duração de 2 horas e serão realizadas todas as quintas-feiras, às 19 horas, na UTFPR, com início em 3 de setembro e término em 26 de novembro de 2026, totalizando 13 encontros.

\subsection*{Conteúdo Programático}

Introdução ao idioma japonês, abordando leitura, escrita, compreensão auditiva e conversação. O conteúdo será desenvolvido por meio de vídeos, textos, explicações gramaticais, vocabulário e atividades práticas de conversação em sala de aula.

\subsection*{Metodologia}

Cada aula será iniciada com uma apresentação do conteúdo previsto. Em seguida, será trabalhado um texto correspondente ao diálogo presente no vídeo da aula, realizando sua leitura, análise do vocabulário e explicação das estruturas linguísticas utilizadas. Posteriormente, os participantes assistirão ao vídeo e aplicarão o vocabulário aprendido em atividades de conversação e perguntas realizadas em sala. Como atividade complementar, serão propostos exercícios para praticar a escrita dos alfabetos \textit{hiragana} e \textit{katakana}, além da introdução gradual aos \textit{kanji}.

\subsection*{Critérios de Avaliação}

Os participantes serão avaliados de forma contínua, considerando o interesse, a participação nas atividades propostas, o envolvimento durante as aulas e a evolução no processo de aprendizagem.

\subsection*{Instrutores e Equipe de Apoio}

\begin{table}[H]
\begin{tabular}{lllll}
\multicolumn{1}{c}{\cellcolor[HTML]{E0E0E0}Nome} & \multicolumn{1}{c}{\cellcolor[HTML]{E0E0E0}CPF} & \multicolumn{1}{c}{\cellcolor[HTML]{E0E0E0}E-mail} & \multicolumn{1}{c}{\cellcolor[HTML]{E0E0E0}Telefone} &  \\
André Luiz da Silva Júnior                       & 084.169.609-84                                  & and695842@gmail.com                                & (41) 9 8830-8688                                     &  \\
\end{tabular}
\end{table}

\end{document}
